% Part: first-order-logic
% Chapter: completeness
% Section: compactness-direct

\documentclass[../../../include/open-logic-section]{subfiles}

\begin{document}

\iftag{FOL}
      {\olfileid{fol}{com}{cpd}}
      {\olfileid{pl}{com}{cpd}}

\olsection{Bukti Langsung Teorema Kekompakan}

Kita bisa mbuktekake Teorema Kekompakan kanthi langsung, tanpa
nggunakake Teorema Kelengkapan, kanthi gagasan kang padha karo bukti
teorema kelengkapan. Ing bukti Teorema Kelengkapan, kita miwiti saka
himpunan~$\Gamma$ kang konsisten saka !!{sentence}s, nggedhekake dadi
himpunan kang konsisten\iftag{FOL}{, jenuh,}{} lan !!{complete}, yaiku
himpunan~$\Gamma^*$ saka !!{sentence}s, banjur nuduhake yen ing
\iftag{FOL}{model term~$\Struct{M(\Gamma^*)}$}{!!{valuation}~
  $\pAssign{v(\Gamma^*)}$}
kang dibangun saka $\Gamma^*$, kabeh !!{sentence}s saka~$\Gamma$ iku
bener, mula $\Gamma$ bisa dicukupi.

Kita bisa nggunakake metodhe kang padha kanggo nuduhake yen himpunan
ukara kang bisa dicukupi kanthi winates uga bisa dicukupi. Kita mung
kudu mbuktekake versi asil-asil kang cocog lan tumuju marang lemma
bebener, kanthi ngganti “konsisten” dadi “bisa dicukupi kanthi
winates.”

\begin{prop}
\ollabel{prop:fsat-ccs}
Upamana $\Gamma$ iku !!{complete} lan bisa dicukupi kanthi winates.
Banjur:
\begin{enumerate}
\tagitem{prvAnd}{$(!A \land !B) \in \Gamma$
  yen lan mung yen $!A \in \Gamma$ lan $!B \in \Gamma$ loro-lorone.}{}

\tagitem{prvOr}{$(!A \lor !B) \in \Gamma$ yen lan mung yen
  $!A \in \Gamma$ utawa $!B \in \Gamma$.}{}

\tagitem{prvIf}{$(!A \lif !B) \in \Gamma$ yen lan mung yen
  $!A \notin \Gamma$ utawa $!B \in \Gamma$.}{}
\end{enumerate}
\end{prop}

\tagprob{FOL}
\begin{prob}
Buktekna \olref[fol][com][cpd]{prop:fsat-ccs}. Aja nggunakake $\Proves$.
\end{prob}
\tagendprob

\tagprob{notFOL}
\begin{prob}
Buktekna \olref[pl][com][cpd]{prop:fsat-ccs}. Aja nggunakake $\Proves$.
\end{prob}
\tagendprob

\iftag{FOL}{%
\begin{lem}
\ollabel{lem:fsat-henkin} Saben himpunan~$\Gamma$ kang bisa dicukupi
kanthi winates bisa digedhekake dadi himpunan jenuh~$\Gamma'$ kang
bisa dicukupi kanthi winates.
\end{lem}
}{}

\tagprob{FOL}
\begin{prob}
Buktekna \olref[fol][com][cpd]{lem:fsat-henkin}. (Pituduh: Langkah
wigatine yaiku nuduhake yen $\Gamma_n$ bisa dicukupi kanthi winates,
mula $\Gamma_n \cup \{!D_n\}$ uga bisa dicukupi kanthi winates, tanpa
nggunakake !!{derivation}s utawa konsistensi.)
\end{prob}
\tagendprob

\iftag{FOL}{
\begin{prop}\ollabel{prop:fsat-instances}
Upamana $\Gamma$ jangkep, bisa dicukupi kanthi winates, lan jenuh.
\begin{tagenumerate}{prvEx,prvAll}
\tagitem{prvEx}{$\lexists[x][!A(x)] \in \Gamma$ yen lan mung yen
  $!A(t) \in \Gamma$ kanggo paling ora siji term katutup~$t$.}{}
\tagitem{prvAll}{$\lforall[x][!A(x)] \in \Gamma$ yen lan mung yen
  $!A(t) \in \Gamma$ kanggo kabeh term katutup~$t$.}{}
\end{tagenumerate}
\end{prop}
}{}

\tagprob{FOL}
\begin{prob}
Buktekna \olref[fol][com][cpd]{prop:fsat-instances}.
\end{prob}
\tagendprob

\begin{lem}
\ollabel{lem:fsat-lindenbaum} Saben himpunan~$\Gamma$ kang bisa
dicukupi kanthi winates bisa digedhekake dadi himpunan~$\Gamma^*$
kang !!a{complete} lan bisa dicukupi kanthi winates.
\end{lem}

\tagprob{FOL}
\begin{prob}
Buktekna \olref[fol][com][cpd]{lem:fsat-lindenbaum}. (Pituduh: langkah
wigatine yaiku nuduhake yen $\Gamma_n$ bisa dicukupi kanthi winates,
mula $\Gamma_n \cup \{!A_n\}$ utawa $\Gamma_n \cup \{\lnot !A_n\}$
bisa dicukupi kanthi winates.)
\end{prob}
\tagendprob

\tagprob{notFOL}
\begin{prob}
Buktekna \olref[pl][com][cpd]{lem:fsat-lindenbaum}. (Pituduh: langkah
wigatine yaiku nuduhake yen $\Gamma_n$ bisa dicukupi kanthi winates,
mula $\Gamma_n \cup \{!A_n\}$ utawa $\Gamma_n \cup \{\lnot !A_n\}$
bisa dicukupi kanthi winates.)
\end{prob}
\tagendprob

\begin{thm}[Kekompakan]
\ollabel{thm:compactness-direct} $\Gamma$ bisa dicukupi yen lan mung
yen bisa dicukupi kanthi winates.
\end{thm}

\begin{proof}
Yen $\Gamma$ bisa dicukupi, mula ana
\iftag{FOL}{!!a{structure}~$\Struct{M}$}{!!a{valuation}~$\pAssign{v}$}
nganti $\iftag{FOL}{\Sat{M}}{\pSat{v}}{!A}$ kanggo kabeh $!A \in
\Gamma$. Mesthi wae, \iftag{FOL}{$\Struct{M}$}{$\pAssign{v}$} iki uga
nyukupi saben subhimpunan winates saka~$\Gamma$, mula $\Gamma$ bisa
dicukupi kanthi winates.

Saiki upamana $\Gamma$ bisa dicukupi kanthi winates.\iftag{FOL}{
Miturut \olref{lem:fsat-henkin}, ana himpunan $\Gamma' \supseteq
\Gamma$ kang bisa dicukupi kanthi winates lan jenuh.}{} Miturut
\olref{lem:fsat-lindenbaum}, $\iftag{FOL}{\Gamma'}{\Gamma}$ bisa
digedhekake dadi himpunan~$\Gamma^*$ kang !!a{complete} lan bisa
dicukupi kanthi winates\iftag{FOL}{, lan $\Gamma^*$ isih jenuh}{}.
Bangunen \iftag{FOL}{model term~$\Struct{M(\Gamma^*)}$}{!!{valuation}~
$\pAssign{v(\Gamma^*)}$} kaya ing
\olref[mod]{defn:termmodel}.\iftag{FOL}{ Catheten yen
\olref[mod]{prop:quant-termmodel} ora gumantung marang kasunyatan yen
$\Gamma^*$ konsisten (utawa !!{complete} utawa jenuh), nanging mung
marang kasunyatan yen $\Struct{M(\Gamma^*)}$ katutup.}{}
Bukti Lemma Bebener (\olref[mod]{lem:truth}) tetep lumaku yen kita
ngganti rujukan menyang \olref[ccs]{prop:ccs}\iftag{FOL}{ lan
\olref[hen]{prop:saturated-instances} nganggo rujukan menyang
\olref{prop:fsat-ccs} lan \olref{prop:fsat-instances}}{ nganggo
pratelan babagan himpunan jangkep kang bisa dicukupi kanthi winates
ing ndhuwur}. Cathetan koreksi sumber OLPL-055: cabang logika
proposisional ing sumber mungkasi ukara sawise rujukan lawas; teks
iki mulihake pangganti prop:fsat-ccs kang dibutuhake cabang kasebut.
\end{proof}

\tagprob{FOL}
\begin{prob}
Tulisen bukti jangkep Lemma Bebener
(\olref[fol][com][mod]{lem:truth}) ing versi kang dibutuhake kanggo
bukti \olref[fol][com][cpd]{thm:compactness-direct}.
\end{prob}
\tagendprob

\tagprob{notFOL}
\begin{prob}
Tulisen bukti jangkep Lemma Bebener
(\olref[pl][com][mod]{lem:truth}) ing versi kang dibutuhake kanggo
bukti \olref[pl][com][cpd]{thm:compactness-direct}.
\end{prob}
\tagendprob

\end{document}
